% chapters/requirements/rf_ingestion.tex
% RF-9: Ingesta de documentos y OCR
% ============================================================================

\item \textbf{Ingesta de documentos y \acs{ocr}.} \label{rf:ingestion}
La ingesta de exámenes se realiza mediante un \dfn{watcher} que monitoriza carpetas
\acs{sftp} en busca de nuevos ficheros \acs{pdf}. Cada \acs{pdf} corresponde a un único
alumno. Tras la recepción, se encola una tarea de \acs{ocr} que
extrae los atributos definidos en las zonas de reconocimiento del modelo correspondiente
(nombre, \acs{dni}, \acs{nia}, etc.) y se crea automáticamente la instancia de examen. El
motor de \acs{ocr} se diseña como un \dfn{plugin} intercambiable, permitiendo mantener
varios motores y seleccionar cuál usar de forma dinámica.

\begin{functional}

    % RF-9.1 ----------------------------------------------------------------
    \item \textbf{Monitorización de carpetas \acs{sftp}.} \label{rf:ingestion:watcher}
    El sistema debe mantener un \dfn{watcher} que monitorice una o varias carpetas
    \acs{sftp} en busca de nuevos ficheros \acs{pdf}.

    \textbf{Entrada:} aparición de un nuevo fichero \acs{pdf} en cualquiera de las carpetas
    \acs{sftp} configuradas.

    \textbf{Proceso:} el \dfn{watcher} detecta la presencia del nuevo fichero, lo valida
    como \acs{pdf} válido y lo transfiere al almacenamiento temporal para su procesamiento.
    Debe soportar múltiples carpetas simultáneamente (una por organización o por examen). El
    \dfn{watcher} se ejecuta como un proceso demonio independiente.

    \textbf{Salida:} el fichero queda disponible para la siguiente fase de ingesta. Se
    genera un evento de recepción en el \textit{log} de aplicación.

    % RF-9.2 ----------------------------------------------------------------
    \item \textbf{Ingesta de \acs{pdf}.} \label{rf:ingestion:ingest}
    El sistema debe procesar cada \acs{pdf} recibido: validarlo, almacenarlo en MinIO y
    encolar su tarea de \acs{ocr}.

    \textbf{Entrada:} evento de recepción de un nuevo \acs{pdf} procedente del
    \dfn{watcher}.

    \textbf{Proceso:} el sistema valida que el fichero es un \acs{pdf} válido y no excede
    el tamaño máximo configurado por organización (por defecto, 50~MB). El \acs{pdf} se
    almacena en MinIO (en el \dfn{bucket} de la organización correspondiente) y se crea un
    registro de instancia en estado \texttt{RECEIVED}. A continuación, se encola una tarea
    en la cola de trabajos para el procesamiento \acs{ocr}, pasando la instancia al estado
    \texttt{QUEUED}.

    \textbf{Salida:} instancia creada en estado \texttt{QUEUED} con el \acs{pdf}
    almacenado. Si el fichero no es válido, se registra el error y se descarta.

    % RF-9.3 ----------------------------------------------------------------
    \item \textbf{Asociación de fuente de ingesta a examen.}
    \label{rf:ingestion:source}
    El sistema debe permitir asociar una fuente de ingesta (carpeta \acs{sftp}) a un examen
    concreto, de forma que los \acs{pdf} recibidos se vinculen automáticamente al examen
    correcto.

    \textbf{Entrada:} configuración definida por el coordinador al crear o modificar un
    examen.

    \textbf{Proceso:} la estrategia principal de asociación se basa en la estructura de
    carpetas \acs{sftp}: cada examen se asocia a una carpeta identificada por el
    identificador del examen. Si además se ha definido una zona de reconocimiento para el
    identificador de examen en el modelo, el resultado del \acs{ocr} sobre dicha zona tiene
    prioridad sobre la asociación por carpeta, permitiendo reasignar el \acs{pdf} al examen
    correcto incluso si se depositó en una carpeta errónea.

    \textbf{Salida:} la asociación queda registrada y el \dfn{watcher} la utiliza para
    vincular cada \acs{pdf} al examen correspondiente.

    % RF-9.4 ----------------------------------------------------------------
    \item \textbf{Ejecución de \acs{ocr} sobre zonas definidas.}
    \label{rf:ingestion:ocr_zones}
    El sistema debe ejecutar el motor de \acs{ocr} sobre las zonas de reconocimiento
    definidas en el modelo del examen para extraer los atributos del \acs{pdf}.

    \textbf{Entrada:} tarea de \acs{ocr} desencolada por un \dfn{worker} de Celery, con la
    referencia a la instancia y el modelo asociado.

    \textbf{Proceso:} el \dfn{worker} recupera el \acs{pdf} de MinIO y, para cada zona
    definida en el modelo, recorta la región de la página correspondiente y la envía al
    motor de \acs{ocr} seleccionado. Los resultados se almacenan como atributos de la
    instancia. Si un atributo no puede extraerse, se activa el \textit{flag}
    \texttt{has\_issues}. La instancia pasa al estado \texttt{PENDING\_GRADING} tras
    finalizar el procesamiento.

    \textbf{Salida:} instancia actualizada con los atributos extraídos.

    % RF-9.5 ----------------------------------------------------------------
    \item \textbf{\textit{Matching} inteligente de atributos.}
    \label{rf:ingestion:matching}
    El sistema debe intentar asociar los textos extraídos por el \acs{ocr} a atributos
    válidos del examen, incluso sin una coincidencia exacta.

    \textbf{Entrada:} texto extraído por el \acs{ocr} para un atributo concreto (nombre,
    \acs{dni}, \acs{nia}).

    \textbf{Proceso:} el sistema compara el texto extraído contra la lista de candidatos
    válidos (alumnos convocados al examen). Para el nombre, se utiliza una métrica de
    similitud textual (distancia de Levenshtein o similitud difusa). Para el \acs{dni} y el
    \acs{nia}, se aplica tolerancia a errores comunes de \acs{ocr} (confusión de caracteres
    similares: 0/O, 1/l/I, etc.). Se establece un umbral de confianza configurable por
    debajo del cual la coincidencia se marca como incierta y se genera incidencia.

    \textbf{Salida:} atributo asignado con un nivel de confianza, o incidencia si no se
    alcanza el umbral.

    % RF-9.6 ----------------------------------------------------------------
    \item \textbf{Selección dinámica de motor \acs{ocr}.}
    \label{rf:ingestion:ocr_select}
    El sistema debe permitir seleccionar el motor de \acs{ocr} a utilizar, tanto de forma
    global (por organización) como por petición individual.

    \textbf{Entrada:} configuración global de la organización (motor por defecto) o
    parámetro opcional en la petición de ejecución de \acs{ocr}.

    \textbf{Proceso:} el sistema consulta la configuración del motor \acs{ocr}: si la
    petición incluye un motor específico, se usa ese; en caso contrario, se usa el motor por
    defecto definido en la configuración de la organización. Si el motor solicitado no está
    disponible o registrado, se rechaza la petición.

    \textbf{Salida:} la tarea de \acs{ocr} se ejecuta con el motor seleccionado.

    % RF-9.7 ----------------------------------------------------------------
    \item \textbf{Registro de motores \acs{ocr} \dfnpl{plugin}.}
    \label{rf:ingestion:ocr_registry}
    El sistema debe permitir registrar y mantener varios motores de \acs{ocr} como
    \dfnpl{plugin} intercambiables.

    \textbf{Entrada:} configuración del sistema que define los motores disponibles (por
    ejemplo, EasyOCR, Tesseract, servicios externos).

    \textbf{Proceso:} cada motor implementa una interfaz definida por el sistema que acepta
    una imagen (recorte de la zona) y devuelve el texto reconocido con un nivel de
    confianza. Los motores se registran en la configuración del sistema con un identificador
    único, un nombre descriptivo y los parámetros propios del motor (idioma, modelo
    entrenado, etc.).

    \textbf{Salida:} el motor queda disponible para ser seleccionado como motor por defecto
    por organización o por petición.

    % RF-9.8 ----------------------------------------------------------------
    \item \textbf{\acs{ocr} sobre anotación de calificación.}
    \label{rf:ingestion:ocr_grade}
    El sistema debe soportar la ejecución de \acs{ocr} sobre una anotación de calificación
    creada por un corrector, extrayendo la nota y asignándola al problema correspondiente.

    \textbf{Entrada:} petición POST al recurso de \acs{ocr} de calificación, con la
    referencia a la anotación que contiene la nota (manuscrita con \textit{stylus} o en
    texto plano) y el identificador del problema al que se refiere.

    \textbf{Proceso:} el sistema extrae el contenido de la anotación (imagen de trazos o
    texto), ejecuta el motor de \acs{ocr} para reconocer el valor numérico y, si se
    reconoce con confianza suficiente, asigna automáticamente la calificación al problema
    indicado. Si la confianza es insuficiente, se devuelve el texto reconocido sin aplicar
    la calificación, permitiendo al corrector confirmar o corregir manualmente.

    \textbf{Salida:} respuesta con el valor reconocido, el nivel de confianza y la
    indicación de si la calificación se asignó automáticamente o queda pendiente de
    confirmación.

    % RF-9.9 ----------------------------------------------------------------
    \item \textbf{\acs{ocr} sobre anotaciones de \textit{stylus} (extensible).}
    \label{rf:ingestion:ocr_stylus}
    El sistema debe soportar la ejecución de \acs{ocr} sobre anotaciones de \textit{stylus}
    para convertir trazos manuscritos a texto.

    \textbf{Entrada:} petición POST al recurso de \acs{ocr} de anotaciones, con la
    referencia a una o varias anotaciones de tipo \textit{stylus}.

    \textbf{Proceso:} el sistema renderiza los trazos \acs{json} de la anotación a una
    imagen y la envía al motor de \acs{ocr} seleccionado. En la versión~1.0 se implementa
    la interfaz y el flujo completo; la calidad del reconocimiento dependerá del motor
    utilizado y podrá mejorarse en versiones futuras.

    \textbf{Salida:} respuesta con el texto reconocido y el nivel de confianza.

\end{functional}
